\section{Modelo: el ciclo de Otto de un qubit}

\begin{frame}{Calor y trabajo en dos niveles}
    \[
        \avg{E} = p\,\omega
        \qquad\Longrightarrow\qquad
        d\avg{E} = \underbrace{\omega\,dp}_{\text{calor}} + \underbrace{p\,d\omega}_{\text{trabajo}}
    \]

    \vspace{4pt}
    \begin{itemize}
        \item \textbf{Calor}: cambia la ocupación, la brecha fija (contacto con un baño).
        \item \textbf{Trabajo}: cambia la brecha, la ocupación fija (cambio rápido de $\omega$).
    \end{itemize}

    \vspace{4pt}
    Ocupación de equilibrio con un baño a temperatura $T$ (Boltzmann):
    \[
        a(\omega, T) = \frac{1}{1+e^{\omega/k_BT}}
    \]

    {\small $\hbar = k_B = 1$;\quad $T_h = 1$,\; $T_c = 0{,}5$ \;$\Rightarrow$\; $\eta_C = 0{,}5$.}
\end{frame}

\begin{frame}{El ciclo de Otto}
    \begin{columns}[c]
        \begin{column}{0.48\textwidth}
            \begin{enumerate}
                \setlength{\itemsep}{3pt}
                \item Contacto caliente ($\omega_h$, $\tau_h$)
                \item Bajar la brecha $\omega_h\to\omega_c$
                \item Contacto frío ($\omega_c$, $\tau_c$)
                \item Subir la brecha $\omega_c\to\omega_h$
            \end{enumerate}
            \vspace{6pt}
            \[
                W = (\omega_h-\omega_c)(p_1-p_2),\qquad Q_h = \omega_h(p_1-p_2)
            \]
            \[
                \eta = \frac{W}{Q_h} = 1-\frac{\omega_c}{\omega_h}
            \]
        \end{column}
        \begin{column}{0.48\textwidth}
            \centering
            \begin{tikzpicture}[scale=1.0, >=stealth]
                \draw[->] (0,0) -- (4.6,0) node[below] {\small $\omega$};
                \draw[->] (0,0) -- (0,3.6) node[left] {\small $p$};
                \coordinate (A) at (3.6,0.9);
                \coordinate (B) at (3.6,2.8);
                \coordinate (C) at (1.4,2.8);
                \coordinate (D) at (1.4,0.9);
                \draw[very thick, usmRed, ->] (A) -- node[right] {\scriptsize 1: $Q_h$} (B);
                \draw[very thick, usmBlue, ->] (B) -- node[above] {\scriptsize 2: trabajo} (C);
                \draw[very thick, usmGreen, ->] (C) -- node[left] {\scriptsize 3: $Q_c$} (D);
                \draw[very thick, usmBlue, ->] (D) -- node[below] {\scriptsize 4: trabajo} (A);
                \draw[dashed, gray] (B) -- (0,2.8) node[left] {\small $p_1$};
                \draw[dashed, gray] (A) -- (0,0.9) node[left] {\small $p_2$};
                \draw[dashed, gray] (D) -- (1.4,0) node[below] {\small $\omega_c$};
                \draw[dashed, gray] (A) -- (3.6,0) node[below] {\small $\omega_h$};
                \foreach \P in {A,B,C,D} \fill (\P) circle (1.5pt);
            \end{tikzpicture}
        \end{column}
    \end{columns}
\end{frame}

\begin{frame}{Tiempo finito}
    Relajación hacia el equilibrio en cada contacto:
    \[
        p(\tau) = a + (p_0 - a)\,e^{-\gamma\tau}, \qquad x \equiv e^{-\gamma \tau}
    \]
    Ciclo límite y potencia:
    \[
        p_1 - p_2 = (a_h - a_c)\,\frac{(1-x_h)(1-x_c)}{1 - x_h x_c},
        \qquad
        \avg{P} = \frac{(\omega_h-\omega_c)(p_1-p_2)}{\tau_h + \tau_c + 2\tau_q}
    \]
    \begin{itemize}
        \item Contacto corto: poco calor. Contacto largo: ciclo lento.
        \item \textbf{Carnot solo con potencia cero.}
    \end{itemize}
\end{frame}

\begin{frame}{El trabajo fluctúa}
    Trabajo por ciclo, con $s\in\{0,1\}$ el estado del qubit:
    \[
        w = (\omega_h - \omega_c)(s_1 - s_2)
    \]
    Los ciclos forman una cadena de Markov. Matriz inclinada:
    \[
        \theta(\chi) = \ln \lambda_{\max}\!\big[M_h\, E(\chi)\, M_c\, E(-\chi)\big],
        \qquad E(\chi) = \operatorname{diag}\!\big(1, e^{\chi(\omega_h-\omega_c)}\big)
    \]
    $\theta$ = generadora de cumulantes (análoga a una energía libre):
    \[
        \avg{P} = \frac{\theta'(0)}{\tau_{\text{ciclo}}},\qquad
        \dP = \lim_{t\to\infty}\frac{\operatorname{Var}\,W(t)}{t} = \frac{\theta''(0)}{\tau_{\text{ciclo}}}
    \]
    {\small $\det$ no depende de $\chi$ $\Rightarrow$ forma cerrada. Validado con Monte Carlo.}
\end{frame}
